\begingroup
\setlength{\LTcapwidth}{\textwidth}
\small
\setlength{\tabcolsep}{4pt}
\captionsetup{font=normalsize}
\begin{longtable}{@{}p{2.4cm}p{2.8cm}>{\raggedright\arraybackslash}p{2.558cm}>{\raggedright\arraybackslash}p{2.558cm}>{\raggedright\arraybackslash}p{2.558cm}@{}}
\caption{Ringkasan penelitian terkait pada prediksi glukosa,
    pengetahuan klinis, dan retrieval.}
\label{tab:tab2.5_penelitian_terkait} \\
\toprule
\textbf{Peneliti} &
        \textbf{Fokus / Pendekatan} &
        \textbf{Jenis Data} &
        \textbf{Temuan Utama} &
        \textbf{Keterbatasan Relevan} \\
\midrule
\endfirsthead

\multicolumn{5}{@{}l}{\small\emph{Tabel \thetable{} (lanjutan)}} \\
\toprule
\textbf{Peneliti} &
        \textbf{Fokus / Pendekatan} &
        \textbf{Jenis Data} &
        \textbf{Temuan Utama} &
        \textbf{Keterbatasan Relevan} \\
\midrule
\endhead

\bottomrule
\endlastfoot

Woldaregay dkk. (2019) &
Prediksi glukosa berbasis data &
Data pemantauan dan variabel pendukung &
Pemodelan berbasis data layak untuk prediksi dinamika glukosa &
Fokus utama pada prediksi, belum pada retrieval pengetahuan klinis \\

Ghimire dkk. (2024) &
Generalisasi model prediksi glukosa &
Data dari karakteristik/dataset berbeda &
Performa dapat berubah pada distribusi data yang berbeda &
Generalisasi menjadi tantangan untuk penerapan lintas kondisi \\

Sarani Rad dkk. (2024) &
Personal health knowledge graph &
Data pasien dan pengetahuan terstruktur &
Integrasi pengetahuan dapat mendukung layanan personal &
Membutuhkan representasi pengetahuan dan integrasi yang lebih kompleks \\

Kresevic dkk. (2024) &
RAG untuk penafsiran pedoman klinis &
Dokumen klinis dan konteks klinis &
Retrieval dapat menyediakan konteks pedoman untuk pembangkitan &
Kueri didasarkan pada kondisi yang tersedia saat peninjauan \\

Lee dkk. (2024) &
RAG berbasis pedoman diabetes &
Dokumen pedoman diabetes dan pertanyaan &
Grounding pada pedoman dapat membantu mengurangi halusinasi &
Kueri bersumber dari pertanyaan pengguna \\

Wang dkk. (2024) &
RAG untuk edukasi diabetes &
Dokumen klinis dan pertanyaan pengguna &
Keluaran dapat dinilai dari akurasi, kelengkapan, dan keterpahaman &
Retrieval dipicu oleh pertanyaan pengguna \\

Jiang dkk. (2023) &
Active retrieval augmented generation &
Informasi yang dibutuhkan selama generation &
Retrieval dapat dipicu secara dinamis selama pembangkitan &
Conditioning berasal dari kebutuhan model bahasa, bukan prediktor fisiologis eksternal \\
\end{longtable}
\endgroup
